%%%
\chapter{Introduction}
%%%
\section{Introduction}

Mobile, semi-connected computing devices have quickly become one of the most dominant forms of personal data storage, data retrieval, and communications  devices on the planet.  Convergence has blurred the lines between what is a computation device, data storage/retrieval device, entertainment device, and communications device, to the point where many mobile devices have features of all four device types in them.  Mobile phones in particular can provide telephony, e-mail, instant messaging, music and video playback, games, imaging, and personal/business databases.  Such devices can run sophisticated software packages, and can often be used in the place of bulkier personal computer devices such as laptop computers.

Unfortunately, we still live in a world where network availability is not completely ubiquitous.  Cellular devices often have issues connecting to their network inside buildings, under overpasses, inside tunnels, and in sparsely populated locales.  For these reasons, applications written for mobile devices, particularly those which rely on some form of database, typically continue to require local storage to ensure data availability when network service is inaccessible.  Keeping this information up-to-date with respect to other databases requires data synchronization.

Of particular note, data synchronization for most mobile devices is still fairly primitive.  It may require the user to plug their mobile device in to a more traditional personal computer style device via a cable, and even when a wireless connection for synchronization is available, the user usually must initiate the synchronization process manually.  The device is typically unusable until the synchronization process has completed.  This may be acceptable for small databases which don�t change frequently (such as a personal address book), however for more complicated applications in business, government, and medicine, the requirements to be physically present with the device to be synchronized against (or potentially through), to manually enable and initialize the synchronization process, and to wait for the process to complete may be more onerous, and may work against the desire to ensure that the synchronized data is as up-to-date as is reasonably possible at the time the data is required.  For large data sets which change frequently, these requirements can become problematic, and can result in the user making decisions based on invalid data, and a reduced desire to use the tool for this purpose. Lengthy and onerous synchronization procedures may result in the user becoming lax in ensuring they synchronize frequently, increasing the chances that a cost may be incurred based on using invalid data.

Fortunately, mobile devices are becoming sufficiently powerful such that they can employ intelligence coupled with automated background synchronization to avoid these issues.  However, such intelligent synchronization systems have by and large not been put into place in favour of the status quo.  

Another recent innovation in many mobile devices is the ability to connect to multiple networks.  Mobile phones such as the Apple iPhone and Palm Treo are capable of connecting to cellular (GPRS, EDGE) and WiFi (802.11.a/b/g) networks for data transfer.  Each of these different network types typically have completely different data transfer rates and costs associated with them, with even identical networking technologies having different costs at different locales.  These differences can also affect the cost associated with data synchronization.

This project will propose an improved mechanism for synchronizing data with such devices, named Expedient Trickle Sync (ETS).  ETS is an intelligent, automated data synchronization system that synchronizes in the background, with the goal of reducing the overall cost of data synchronization in the following two areas:
\begin{enumerate}
\item By reducing the real cost associated with network access, and
\item By reducing the virtual cost associated with the users use of invalid or out-of-date data.
\end{enumerate}

\section{Research Scope}
The scope of this research is defined to be:
\begin{enumerate}
\item Identify and precisely specify the goals for the ETS algorithm,
\item Implement the ETS algorithm in a computer programming language,
\item Implement a simulation environment for evaluating the cost of a data synchronization algorithm,
\item Evaluate the ETS algorithm by running it through the simulator, along with both a purely random algorithm and a common, existing industry algorithm (HotSync).
\end{enumerate}

\section{Methodology}
Development of this project will be undertaken in three phases: the development of a synchronization algorithm simulation environment, the development and definition of the Expedient Trickle Sync algorithm, and the evaluation of the ETS algorithm within the simulation environment.

\subsection{Simulator}
Phase one will see the development of a simulator environment capable of running a given synchronization strategy algorithm across the course of many simulated days, calculating along the way the cost of synchronization using the strategy.  After a set number of simulated days, an average cost will be calculated and stored for the given sync algorithm.

The simulation will simulate a variety of elements, including the movement of the user from location to location during their simulated day (based on a definable schedule), the networks available in each location (and their associated cost , speed, and probability of signal loss), the users changing activity level in relation to the semi-connected mobile device, changes to the host database (which the mobile device synchronizes against), and user record access.  These latter two factors will be coded to occur based on Poisson processes \cite{RefWorks:32} (for determining when modifications or accesses are to occur), and a Normal distribution (for determining which records are affected by the modifications and accesses ascertained by the Poisson process).

\subsection{ETS Algorithm}
Phase two of the project is the definition of the ETS algorithm. This algorithm uses pattern analysis and intelligent heuristics to accomplish the following:
\begin{enumerate}
\item Make intelligent decisions as to when to initiate background synchronization sessions.  The decision as to when to synchronize is to be based on previous access patterns, with an attempt to minimize the real cost by synchronizing as much data as possible when in contact with low-cost, high-speed networks, and as little as possible when in contact with high-cost, low speed networks.  At the same time, it will attempt to minimize the virtual cost associated with using expired data by attempting to ensure that synchronization is frequent enough to maximize the probability that at time t, a record R required by the user is up-to-date, and
\item During each synchronization session, order the records to be synchronized so as to prioritize those records which have a high probability of being required between the time of the current synchronization ($t_{curr}$) and the next synchronization session ($t_{next}$), potentially synchronizing only a subset of the available modified records based on their predicted likelihood of being required during this time, if the network cost at $t_{curr}$ is significant.
\end{enumerate}

As the user is presumed to be mobile, the algorithm will need to deal gracefully with unexpected disconnection.  As the algorithm is a learning algorithm, it is expected to begin with no knowledge as to the users patterns or data needs, and will thus record these issues during the course of each day, improving its heuristics.

\subsection{Usage Model Study}
In order to obtain a real-world model 

\subsection{Evaluation}
The third and final phase will use the simulator to evaluate three algorithms against our simulated user:  the ETS algorithm, an algorithm analogous to the Palm HotSync �Slow Sync� protocol \cite{RefWorks:32, RefWorks:24, RefWorks:31}, an algorithm analogous to the Palm HotSync �Fast Sync� protocol \cite{RefWorks:32, RefWorks:24, RefWorks:31}, a random algorithm (which synchronizes random data at random times, used as a control), and simply not synchronizing at all (the "null synchronizer", also used as a control).  Each algorithm will be executed for a large number of simulated days, based on a model user, after which an average cost per day will be derived.  This average cost will be expressed as a simple polynomial, with a partial polynomial solution being provided through simplification using real-world cost profiles for data access on the various simulated network types.  These simplified polynomials will then be compared to evaluate how the ETS algorithm compares to the other two in terms of cost savings for both the real network access cost, and the virtual cost associated with using out-of-date or invalid data.  The simulation will be re-run across a series of model users, to evaluate how each algorithm performs under different usage scenarios.

The user models will be generated based on a real-world data, obtained from a proposed study of existing organizations which use existing synchronization solutions.  Details on the proposed study can be found in \ref{proposed_study}.

\subsection{Target Platform}
The language used for developing and deploying the simulator the simulator will be Objective-C v2.0 \cite{RefWorks:35} on Apple Inc. hardware.  While initially the simulation application will be developed to run a stand-alone single day simulation, if the simulation time is sufficiently large, the application may be modified to use an Xgrid controller \cite{RefWorks:36} to farm out simulated days to multiple systems running in parallel.
